% Integration note:
% This file is designed as an integration-ready frontier synthesis section.
% It assumes the manuscript preamble has already defined theorem environments
% and the claim-status macros such as \ClaimStatusProvedHere and
% \ClaimStatusOpen.

\section{Higher-genus anomaly classes and the first quantum equation in celestial BF}
\label{sec:celestial-bf-frontier-synthesis}

The governing mechanism of this section is the same at two scales.
Globally, higher-genus quantum corrections are organized by a completed
Maurer--Cartan problem in the modular deformation complex.  Locally, the
first computable quantum correction in celestial holomorphic BF appears as
an obstruction to closing the smallest axion-enlarged low-mode shell under
Jacobi.
The first point is controlled by filtered rigidity in pronilpotent dg Lie
algebras; the second by a determinant line on the first positive KK shell.
The payoff is that both the global higher-genus BV/BRST problem and the
local genus-$1$ celestial BF problem reduce to explicit cohomological
classes rather than to an undifferentiated cloud of ``higher operations.''

The section has four parts.
First, we record the theorem-shaped higher-genus BV/BRST target and then
prove a filtered comparison principle which removes the last formal
uniqueness issue once a renormalized anomaly class has been globalized.
Second, we prove the first multichannel trace-detection theorem beyond the
one-channel locus.
Third, we compute the first positive KK shell of celestial BF from first
principles and show that its determinant sector is formally abelian.
Fourth, we adjoin the lowest axionic generator
$\widetilde\beta_{0,0}$ and compute the first actual Jacobi equation for
$D'_{1,1}$.
The outcome is a new exact statement: the low-shell Jacobi defect is a
single universal tensor, and any nonzero one-loop coefficient
$D'_{1,1}$ must be carried by an explicit off-shell return map.

\subsection{First principles: modular Maurer--Cartan classes and renormalized anomalies}
\label{subsec:frontier-renormalized-anomaly}

At genus~$0$, the manuscript already identifies the bar differential with
BRST in the proved range.  The unresolved higher-genus content is more
precise: one must compare two Maurer--Cartan elements in the same completed
modular deformation problem.
On one side sits the universal modular class $\Theta_{\mathcal A}$ already
constructed in the modular Koszul package; on the other sits the
renormalized perturbative BRST anomaly class coming from local higher
operations.  The 2024 paper of Gaiotto--Kulp--Wu provides exactly the local
input: the anomaly brackets are packaged by an odd nilpotent vector
field on the formal space of couplings, and these higher brackets satisfy
quadratic identities interpreted as Wess--Zumino consistency.
The missing step is therefore not to invent a new local algebra, but to
prove that the renormalized globalized anomaly class is the same modular
Maurer--Cartan element as the already-existing universal class.

\begin{theorem}[Higher-genus BV/BRST from the renormalized anomaly class;
\ClaimStatusOpen]
\label{thm:frontier-bv-brst-renormalized-anomaly}
Let $\mathcal A$ be a modular Koszul chiral algebra on a smooth projective
curve $X$, realized as the boundary chiral algebra of a
holomorphic-topological field theory with local perturbative BRST anomaly
brackets in the sense of Gaiotto--Kulp--Wu~\cite{GKW24}.  Assume:
\begin{enumerate}[label=\textup{(\roman*)}]
\item the disk-local perturbative packet agrees with the genus-$0$ local
  bar packet, so that the local perturbative binary and ternary anomaly
  brackets match the compactified Fulton--MacPherson residue operations;
\item the renormalized local Feynman forms extend from smooth
  configuration spaces to logarithmic forms on a hybrid
  Fulton--MacPherson/Deligne--Mumford compactification of stable
  configurations over $\overline{\mathcal M}_{g,n}$, with boundary
  restrictions matching stable-graph factorization;
\item the renormalization scheme is local, factorization-compatible, and
  clutching-compatible, so that boundary counterterms agree with the
  gluing terms in the completed modular deformation complex
  $\operatorname{Def}^{\mathrm{mod}}(\mathcal A)$;
\item the resulting genus-$0$ normalization agrees with the already-proved
  genus-$0$ BRST/bar comparison.
\end{enumerate}
Then there exists a canonical renormalized anomaly class
\[
\alpha^{\mathrm{ren}}_{\mathcal A}
\in
\operatorname{MC}\!\bigl(\operatorname{Def}^{\mathrm{mod}}(\mathcal A)\bigr)
\]
whose image under the modular realization map equals the universal modular
Maurer--Cartan class $\Theta_{\mathcal A}$.
Equivalently,
\[
\alpha^{\mathrm{ren}}_{\mathcal A} \sim \Theta_{\mathcal A}
\]
in the Deligne--Getzler--Hinich Maurer--Cartan $\infty$-groupoid, and on
any one-channel rigid locus this gauge-equivalence upgrades to strict
identity.
Consequently the higher-genus BRST differential is the modular bar
differential twisted by $\Theta_{\mathcal A}$, and the scalar trace of the
renormalized anomaly class recovers the tautological genus package.
\end{theorem}

\begin{remark}[Logical position]
\label{rem:frontier-renormalized-anomaly-position}
Theorem~\ref{thm:frontier-bv-brst-renormalized-anomaly} is not an existence
statement of the same kind as the modular MC construction.
The universal class $\Theta_{\mathcal A}$ is already present.
What remains is a comparison theorem:
\[
\Theta_{\mathcal A} \stackrel{?}{=} \alpha^{\mathrm{ren}}_{\mathcal A}.
\]
Once the local packet is identified, the residual difficulty is global and
modular: extend renormalized amplitudes to stable compactifications,
identify Stokes boundary terms with modular gluing, and then prove that the
resulting Maurer--Cartan lift is the unique one with the prescribed
normalization and trace package.
\end{remark}

\subsubsection*{Proof architecture}

The proof breaks into four steps.

\emph{Step~1: local anomaly brackets $=$ local bar operations.}
Use the higher perturbative brackets of~\cite{GKW24} and identify them with
the compactified local bar residues.
At genus~$0$ this reduces to the binary packet on $C_2$ and the compactified
ternary packet on $\overline M_{0,4}$.

\emph{Step~2: globalize to stable moduli.}
Construct a hybrid compactification of stable configurations over
$\overline{\mathcal M}_{g,n}$ and extend the renormalized Feynman forms to
logarithmic forms on it.

\emph{Step~3: Stokes $=$ modular Maurer--Cartan.}
Push the logarithmic forms to the stable moduli space; fiberwise Stokes then
identifies the sum of boundary terms with the modular differential and the
convolution bracket.
This produces a Maurer--Cartan element
$\alpha^{\mathrm{ren}}_{\mathcal A}$.

\emph{Step~4: compare with the universal class.}
Once $\alpha^{\mathrm{ren}}_{\mathcal A}$ and $\Theta_{\mathcal A}$ live in
the same completed deformation problem, the remaining comparison is a
filtered uniqueness problem.
That filtered uniqueness problem is no longer open: it is the content of the
next proposition.

\subsubsection*{A proved filtered comparison principle}

\begin{proposition}[Filtered trace-detecting uniqueness of modular lifts;
\ClaimStatusProvedHere]
\label{prop:frontier-filtered-trace-detecting-uniqueness}
Let $M$ be a complete descending filtered dg Lie algebra over a field of
characteristic $0$ with filtration
\[
F^1M \supseteq F^2M \supseteq \cdots,
\qquad
M \cong \varprojlim_m M/F^mM,
\]
such that
\[
d(F^mM) \subseteq F^mM,
\qquad
[F^pM,F^qM] \subseteq F^{p+q}M.
\]
Let $V$ be a complete filtered abelian dg Lie algebra, and let
\[
\operatorname{tr}\colon M \longrightarrow V
\]
be a filtration-preserving dg Lie morphism.
Let $\Theta,\Xi\in \operatorname{MC}(M)$, and write
\[
d_{\Theta}:=d+[\Theta,-]
\]
for the twisted differential.
Assume:
\begin{enumerate}[label=\textup{(\roman*)}]
\item $\Theta\equiv \Xi \pmod{F^1M}$;
\item the trace classes agree in $H^1(V)$:
  \[
  [\operatorname{tr}(\Theta)] = [\operatorname{tr}(\Xi)];
  \]
\item for every $m\ge 1$, the induced map
  \[
  H^1\!\bigl(\operatorname{gr}_F^m(M,d_{\Theta})\bigr)
  \xrightarrow{\ \operatorname{gr}^m(\operatorname{tr})\ }
  H^1\!\bigl(\operatorname{gr}_F^mV\bigr)
  \]
  is injective.
\end{enumerate}
Then $\Theta$ and $\Xi$ are gauge-equivalent in the
Deligne--Getzler--Hinich Maurer--Cartan $\infty$-groupoid of~$M$.
\end{proposition}

\begin{proof}
We construct inductively Maurer--Cartan elements
$\Xi^{(m)}\in \operatorname{MC}(M)$, all gauge-equivalent to~$\Xi$, such
that
\[
\Xi^{(m)}-\Theta \in F^{m+1}M
\qquad (m\ge 0).
\]
Set $\Xi^{(0)}:=\Xi$.
Hypothesis~\textup{(i)} gives $\Xi^{(0)}-\Theta\in F^1M$.

Assume inductively that
$\delta^{(m)}:=\Xi^{(m)}-\Theta\in F^{m+1}M$.
Let
\[
\bar\delta^{(m)}\in \operatorname{gr}_F^{m+1}M^1
\]
be its leading term.
Since both $\Xi^{(m)}$ and $\Theta$ satisfy the Maurer--Cartan equation,
subtracting the two equations and reducing modulo $F^{m+2}M$ yields
\[
d_{\Theta}(\bar\delta^{(m)})=0
\qquad
\text{in }\operatorname{gr}_F^{m+1}M.
\]
So $\bar\delta^{(m)}$ defines a class in
$H^1(\operatorname{gr}_F^{m+1}(M,d_{\Theta}))$.

Because $\Xi^{(m)}$ is gauge-equivalent to~$\Xi$, its trace class agrees
with that of~$\Xi$.
By hypothesis~\textup{(ii)} this equals the trace class of~$\Theta$.
Hence the leading graded trace class of $\Xi^{(m)}-\Theta$ vanishes:
\[
[\operatorname{gr}^{m+1}(\operatorname{tr})(\bar\delta^{(m)})]=0
\qquad
\text{in }H^1(\operatorname{gr}_F^{m+1}V).
\]
By hypothesis~\textup{(iii)}, the class of $\bar\delta^{(m)}$ itself must
vanish.
Therefore there exists
$\bar\xi^{(m+1)}\in \operatorname{gr}_F^{m+1}M^0$ such that
\[
d_{\Theta}(\bar\xi^{(m+1)})=\bar\delta^{(m)}.
\]
Choose any lift $\xi^{(m+1)}\in F^{m+1}M^0$ of $\bar\xi^{(m+1)}$ and set
\[
\Xi^{(m+1)}:=e^{-\xi^{(m+1)}}*\Xi^{(m)}.
\]
Because $\xi^{(m+1)}\in F^{m+1}M^0$, the
Baker--Campbell--Hausdorff expansion gives
\[
\Xi^{(m+1)}
=
\Xi^{(m)}-d_{\Xi^{(m)}}\xi^{(m+1)}+O(F^{m+2}M)
=
\Xi^{(m)}-d_{\Theta}\xi^{(m+1)}+O(F^{m+2}M).
\]
Reducing modulo $F^{m+2}M$, the leading term of
$\Xi^{(m+1)}-\Theta$ is
\[
\bar\delta^{(m)}-d_{\Theta}(\bar\xi^{(m+1)})=0.
\]
Hence
\[
\Xi^{(m+1)}-\Theta\in F^{m+2}M.
\]
This closes the induction.

Now form the infinite BCH product
\[
g:=\cdots e^{-\xi^{(3)}}e^{-\xi^{(2)}}e^{-\xi^{(1)}}.
\]
Completeness of the filtration and pronilpotence of the gauge action imply
that this product converges in the gauge group of~$M^0$, and by
construction
\[
g*\Xi=\Theta.
\]
Thus $\Theta$ and $\Xi$ are gauge-equivalent.
\end{proof}

\begin{corollary}[Scalar-saturated modular lift criterion;
\ClaimStatusProvedHere]
\label{cor:frontier-scalar-saturated-modular-lift}
In the setting of
Proposition~\ref{prop:frontier-filtered-trace-detecting-uniqueness},
assume in addition that each graded twisted deformation space
\[
H^1\!\bigl(\operatorname{gr}_F^m(M,d_{\Theta})\bigr)
\]
is at most one-dimensional and that the graded trace map is nonzero on it.
Then the injectivity hypothesis of
Proposition~\ref{prop:frontier-filtered-trace-detecting-uniqueness} holds
automatically.
Consequently the pair
\[
\bigl(\text{genus-$0$ normalization},\ \text{scalar trace}\bigr)
\]
determines the full higher-genus Maurer--Cartan lift uniquely up to gauge.
\end{corollary}

\begin{proof}
A nonzero linear map out of a vector space of dimension at most one is
injective.
Apply
Proposition~\ref{prop:frontier-filtered-trace-detecting-uniqueness}.
\end{proof}

\begin{corollary}[BV/BRST comparison from trace detection;
\ClaimStatusProvedHere]
\label{cor:frontier-bv-brst-comparison-from-trace-detection}
Let $M=\operatorname{Def}^{\mathrm{mod}}(\mathcal A)$ be a complete
filtered dg Lie model of the modular deformation problem of a modular
Koszul chiral algebra $\mathcal A$, filtered by positive genus.
Let
\[
\Theta_{\mathcal A},\ \alpha^{\mathrm{ren}}_{\mathcal A}
\in \operatorname{MC}(M)
\]
be Maurer--Cartan elements such that:
\begin{enumerate}[label=\textup{(\roman*)}]
\item their genus-$0$ restrictions agree;
\item their scalar traces agree in cohomology:
  \[
  [\operatorname{tr}(\alpha^{\mathrm{ren}}_{\mathcal A})]
  =
  [\operatorname{tr}(\Theta_{\mathcal A})];
  \]
\item the graded twisted deformation spaces are detected by the scalar
  trace, for instance by the scalar-saturated hypothesis of
  Corollary~\ref{cor:frontier-scalar-saturated-modular-lift}.
\end{enumerate}
Then
\[
\alpha^{\mathrm{ren}}_{\mathcal A}\sim \Theta_{\mathcal A}
\]
in the Maurer--Cartan $\infty$-groupoid.
Thus, once the renormalized anomaly class has been globalized and its
scalar trace identified, the remaining comparison in
Theorem~\ref{thm:frontier-bv-brst-renormalized-anomaly} is not a separate
open uniqueness problem.
\end{corollary}

\begin{proof}
Apply
Proposition~\ref{prop:frontier-filtered-trace-detecting-uniqueness} to the
genus filtration on $\operatorname{Def}^{\mathrm{mod}}(\mathcal A)$.
\end{proof}

\subsection{Two-channel trace detection beyond the one-channel locus}
\label{subsec:frontier-two-channel-trace-detection}

The filtered comparison principle shifts the frontier from ``uniqueness'' to
``detection.''
The first nontrivial case beyond one-channel rigidity is the two-channel
situation.  The key point is that one does not need a literal product
splitting of the full deformation complex; it suffices to quotient by the
mixed sector and show that the mixed sector has no twisted degree-$1$
cohomology on the associated graded.

\begin{proposition}[Two-channel graded trace detection from mixed-sector
vanishing; \ClaimStatusProvedHere]
\label{prop:frontier-two-channel-graded-trace-detection}
Let $M$ be a filtered dg Lie algebra with descending filtration
\[
F^0M:=M \supseteq F^1M \supseteq F^2M \supseteq \cdots,
\qquad
 d(F^mM)\subseteq F^mM,
\qquad
 [F^pM,F^qM]\subseteq F^{p+q}
\quad (p,q\ge 0),
\]
and let $\Theta\in\operatorname{MC}(M)$.
Suppose there exist filtered dg Lie algebras $M_1,M_2$ and a
filtration-preserving dg Lie morphism
\[
q\colon M\longrightarrow M_1\widehat\oplus M_2
\]
with kernel $C:=\ker q$.
Write
\[
\Theta_i:=q_i(\Theta)\in\operatorname{MC}(M_i),
\qquad i=1,2,
\]
and let
\[
d_{\Theta}=d+[\Theta,-],
\qquad
 d_{\Theta_i}=d+[\Theta_i,-].
\]
Assume that for every $m\ge 1$:
\begin{enumerate}[label=\textup{(\roman*)}]
\item the mixed sector contributes no degree-$1$ class on the graded
  genus-$m$ piece:
  \[
  H^1\!\bigl(\operatorname{gr}_F^m(C,d_{\Theta})\bigr)=0;
  \]
\item there exist filtration-preserving dg Lie morphisms to abelian
  filtered dg Lie algebras
  \[
  \operatorname{tr}_i\colon M_i\longrightarrow V_i,
  \qquad i=1,2,
  \]
  such that the induced maps
  \[
  H^1\!\bigl(\operatorname{gr}_F^m(M_i,d_{\Theta_i})\bigr)
  \xrightarrow{\ \operatorname{gr}^m(\operatorname{tr}_i)\ }
  H^1\!\bigl(\operatorname{gr}_F^mV_i\bigr)
  \]
  are injective.
\end{enumerate}
Then the combined trace map
\[
\operatorname{tr}^{(2)}:=(\operatorname{tr}_1\oplus \operatorname{tr}_2)\circ q
\colon M\longrightarrow V_1\oplus V_2
\]
induces an injective map on every graded twisted deformation space:
\[
H^1\!\bigl(\operatorname{gr}_F^m(M,d_{\Theta})\bigr)
\xrightarrow{\ \operatorname{gr}^m(\operatorname{tr}^{(2)})\ }
H^1\!\bigl(\operatorname{gr}_F^m(V_1\oplus V_2)\bigr).
\]
If in addition
\[
H^2\!\bigl(\operatorname{gr}_F^m(C,d_{\Theta})\bigr)=0
\qquad \text{for all }m\ge 1,
\]
then $q$ induces an isomorphism
\[
H^1\!\bigl(\operatorname{gr}_F^m(M,d_{\Theta})\bigr)
\xrightarrow{\ \sim\ }
H^1\!\bigl(\operatorname{gr}_F^m(M_1,d_{\Theta_1})\bigr)
\oplus
H^1\!\bigl(\operatorname{gr}_F^m(M_2,d_{\Theta_2})\bigr)
\]
for every $m$.
\end{proposition}

\begin{proof}
Because $q$ is a dg Lie morphism and $q(\Theta)=(\Theta_1,\Theta_2)$, it
intertwines the twisted differentials:
\[
q\circ d_{\Theta}=(d_{\Theta_1}\oplus d_{\Theta_2})\circ q.
\]
Since $C=\ker q$ is a dg Lie ideal, $d_{\Theta}$ preserves~$C$.
Passing to the associated graded at filtration level~$m$ yields a short
exact sequence of cochain complexes
\[
0\longrightarrow \operatorname{gr}_F^m(C,d_{\Theta})
\longrightarrow \operatorname{gr}_F^m(M,d_{\Theta})
\xrightarrow{\ \operatorname{gr}^mq\ }
\operatorname{gr}_F^m(M_1,d_{\Theta_1})\oplus
\operatorname{gr}_F^m(M_2,d_{\Theta_2})
\longrightarrow 0.
\]
The associated long exact sequence in cohomology contains
\[
H^1\!\bigl(\operatorname{gr}_F^m(C,d_{\Theta})\bigr)
\longrightarrow
H^1\!\bigl(\operatorname{gr}_F^m(M,d_{\Theta})\bigr)
\xrightarrow{\ q_*\ }
H^1\!\bigl(\operatorname{gr}_F^m(M_1,d_{\Theta_1})\bigr)
\oplus
H^1\!\bigl(\operatorname{gr}_F^m(M_2,d_{\Theta_2})\bigr).
\]
Hypothesis~\textup{(i)} kills the left-hand term, so $q_*$ is injective.
Applying the factorwise traces and using hypothesis~\textup{(ii)} shows that
$\operatorname{gr}^m(\operatorname{tr}^{(2)})_*$ is injective as well.

If in addition
$H^2(\operatorname{gr}_F^m(C,d_{\Theta}))=0$,
then the next term in the long exact sequence vanishes and $q_*$ becomes
surjective as well as injective.
Hence $q_*$ is an isomorphism on $H^1$.
\end{proof}

\begin{corollary}[Split two-channel locus; \ClaimStatusProvedHere]
\label{cor:frontier-split-two-channel-trace-detection}
In the setting of
Proposition~\ref{prop:frontier-two-channel-graded-trace-detection},
suppose that $q$ is an isomorphism, equivalently
\[
M\cong M_1\widehat\oplus M_2
\]
as filtered dg Lie algebras.
Then graded trace detection on $M$ is exactly the product of the factorwise
trace-detection statements.
In particular, on any split two-channel semisimple locus built from two
one-channel trace-rigid factors, the pair of scalar traces
\[
(\kappa_1,\kappa_2)
\]
detects the full graded higher-genus deformation class.
\end{corollary}

\begin{proof}
Take $C=0$ in
Proposition~\ref{prop:frontier-two-channel-graded-trace-detection}.
\end{proof}

\begin{corollary}[Two-channel modular lift rigidity;
\ClaimStatusProvedHere]
\label{cor:frontier-two-channel-modular-lift-rigidity}
Let $M$ be a complete filtered dg Lie algebra and let
\[
\Theta,\Xi\in\operatorname{MC}(M)
\]
be Maurer--Cartan elements with the same genus-$0$ restriction.
Assume that $M$ admits a two-channel quotient as in
Proposition~\ref{prop:frontier-two-channel-graded-trace-detection}, that
for all $m\ge 1$ the mixed sector satisfies
\[
H^1\!\bigl(\operatorname{gr}_F^m(C,d_{\Theta})\bigr)=0,
\]
that the factorwise graded traces detect the diagonal channels, and that
the factorwise trace classes of $\Theta$ and $\Xi$ agree in cohomology.
Then $\Theta$ and $\Xi$ are gauge-equivalent.
Thus on the split two-channel semisimple locus,
\[
\bigl(\text{genus-$0$ normalization},\ \kappa_1,\ \kappa_2\bigr)
\]
determines the full higher-genus lift uniquely up to gauge.
\end{corollary}

\begin{proof}
Combine
Proposition~\ref{prop:frontier-two-channel-graded-trace-detection}
with
Proposition~\ref{prop:frontier-filtered-trace-detecting-uniqueness}.
\end{proof}

\begin{remark}[Where the multichannel frontier now lives]
\label{rem:frontier-two-channel-trace-frontier}
Beyond the one-channel locus, graded trace detection fails only through
cohomology in the mixed sector $C=\ker q$.
Thus the next frontier is not ``more than one channel'' by itself, but the
appearance of genuine mixed twisted classes in
\[
H^1\!\bigl(\operatorname{gr}_F^m(C,d_{\Theta})\bigr)
\]
and, if one wants a diagonal normal form rather than mere trace detection,
in
\[
H^2\!\bigl(\operatorname{gr}_F^m(C,d_{\Theta})\bigr).
\]
The split semisimple two-channel case is the first proved multichannel
instance, and every genuinely new non-scalar example must pass through a
nonzero mixed-sector cohomology class.
\end{remark}

\subsection{The first positive KK shell in celestial BF}
\label{subsec:frontier-determinantal-shell}

We now turn from the global modular problem to the first explicit local
quantum shell in celestial holomorphic BF.
The input is Zeng's KK reduction of holomorphic BF theory: the effective
boundary algebra is controlled by the transferred $A_\infty$ operations on
$H_b^{0,\bullet}(S^3)$, with the cubic operation
\[
m_3(a,b,c)=pM(a,hM(b,c))-pM(hM(a,b),c),
\]
the antisymmetry relation
\[
(m_3)_{p,q;r,s}^{u_1,v_1;u_2,v_2}
=
-(m_3)_{r,s;p,q}^{u_2,v_2;u_1,v_1},
\]
and the full boundary OPE formulas for $BB$, $B\widetilde B$, and the
Kodaira--Spencer couplings~\cite{Zeng23}.  The first positive shell already
contains nontrivial higher structure, but it is smaller and more rigid than
one might expect: its entire quantum ambiguity is concentrated on a single
determinant line.

Write
\[
u_0:=b_{0,0},
\qquad
u_1:=b_{1,0},
\qquad
u_2:=b_{0,1},
\qquad
\widetilde\nu_0:=\widetilde b_{0,0}.
\]
Set
\[
U_+:=\operatorname{Span}\{\nu_1,\nu_2\},
\qquad
V_0:=\operatorname{Span}\{\nu_0\},
\qquad
\widetilde V_0:=\operatorname{Span}\{\widetilde\nu_0\}.
\]
Let
\[
\omega\in \Lambda^2U_+^\vee,
\qquad
\omega(\nu_1,\nu_2)=1.
\]

\begin{theorem}[Lowest-mode determinantal rigidity in celestial BF;
\ClaimStatusProvedHere]
\label{thm:frontier-lowest-mode-determinantal-rigidity}
After projection to the lowest KK truncation generated by
$\nu_0,\nu_1,\nu_2,\widetilde\nu_0$, the following hold.
\begin{enumerate}[label=\textup{(\roman*)}]
\item The projected pole-$1$ $B$-valued tree bracket on
  $U_+\times U_+$ vanishes.
\item The only nonzero higher tree-level term is the cubic
  $B\widetilde B$-channel, and it factors through the determinant line:
  \[
  \mu(\nu_i,\nu_j,\widetilde\nu_0)
  =
  -\frac12\,\omega(\nu_i,\nu_j)\,\nu_0.
  \]
  Equivalently,
  \[
  m_{3,\mathrm{low}}(\nu_1,\nu_2,\widetilde\nu_0)=-\frac12\nu_0,
  \qquad
  m_{3,\mathrm{low}}(\nu_2,\nu_1,\widetilde\nu_0)=+\frac12\nu_0.
  \]
\item Every projected $B$-valued bilinear correction of any genus supported
  on $U_+\times U_+\to V_0$ is a unique scalar multiple of
  \[
  \nu(u,v):=\omega(u,v)\,\nu_0.
  \]
\end{enumerate}
In particular, the one-loop adjoint channel on the lowest shell is
necessarily of the form
\[
\delta^{(1)}_{\mathrm{ad,low}}(u,v)
=
\lambda\,h^\vee\,\omega(u,v)\,\nu_0
\]
for a unique scalar $\lambda$, and Zeng's one-loop formula identifies this
scalar with $D'_{1,1}$.
\end{theorem}

\begin{proof}
For the first claim, use the full $BB$ boundary OPE of Zeng.
The $n=2$ term is the pole-$1$ $B$-valued part, and its KK output lies in
mode $(p+r,q+s)$.
For inputs in $U_+$ the only possibilities are
$(2,0)$, $(1,1)$, and $(0,2)$, all outside the projected shell
$U_+\oplus V_0$.
Hence the projected pole-$1$ $B$-valued bracket on $U_+\times U_+$ is zero.

For the cubic term, use the transferred formula
\[
m_3(a,b,c)=pM(a,hM(b,c))-pM(hM(a,b),c).
\]
On $H_b^{0,0}(S^3)\otimes H_b^{0,0}(S^3)\otimes H_b^{0,1}(S^3)$ the second
term vanishes because $h=0$ on $H_b^{0,0}(S^3)$.
Identify
\[
\nu_1\leftrightarrow w_1,
\qquad
\nu_2\leftrightarrow w_2,
\qquad
\widetilde\nu_0\leftrightarrow \epsilon.
\]
Using
\[
h(w_1\epsilon)=\bar w_2,
\qquad
h(w_2\epsilon)=-\bar w_1,
\]
we obtain
\[
m_3(w_1,w_2,\epsilon)=p\bigl(w_1(-\bar w_1)\bigr),
\qquad
m_3(w_2,w_1,\epsilon)=p\bigl(w_2\bar w_2\bigr).
\]
On $S^3$ one has $|w_1|^2+|w_2|^2=1$, and the scalar projector $p$ treats
$|w_1|^2$ and $|w_2|^2$ symmetrically, so each has scalar part $1/2$.
Therefore
\[
m_3(w_1,w_2,\epsilon)=-\frac12,
\qquad
m_3(w_2,w_1,\epsilon)=+\frac12.
\]
This also matches the antisymmetry relation for the structure constants of
$m_3$.

For the third claim, note that
\[
\dim \Lambda^2U_+=1,
\qquad
\dim V_0=1.
\]
Hence
\[
\operatorname{Hom}(\Lambda^2U_+,V_0)\cong \mathbb C,
\]
generated by $\omega\otimes \nu_0$.
Every projected $B$-valued bilinear correction on the first positive shell
is therefore a unique scalar multiple of $\nu$.
Zeng's one-loop adjoint formula for $B^a[p,q]B^b[r,s]$ has coefficient
$(ps-qr)h^\vee f_{ab}{}^c$ and shifted output mode
$(p+r-1,q+s-1)$.
On the pair $(1,0),(0,1)$ this output is exactly $(0,0)$, and on every
other pair in the lowest shell it misses $V_0$.
So the one-loop adjoint term is necessarily of the displayed form, with
$\lambda=D'_{1,1}$.
\end{proof}

\begin{proposition}[Positive-shell determinantal extinction;
\ClaimStatusProvedHere]
\label{prop:frontier-positive-shell-determinantal-extinction}
Let $\mu$ denote the cubic low-shell map of
Theorem~\ref{thm:frontier-lowest-mode-determinantal-rigidity}, and let
$\nu_g\in \operatorname{Hom}(\Lambda^2U_+,V_0)$ be any family of bilinear
corrections, one for each positive genus.
Then every partial operadic insertion among the maps $\mu$ and $\nu_g$
vanishes:
\[
\mu\circ \mu = 0,
\qquad
\mu\circ \nu_g = 0,
\qquad
\nu_g\circ \mu = 0,
\qquad
\nu_g\circ \nu_h = 0.
\]
Equivalently, the determinant sector on the first positive shell is an
abelian graded Lie algebra of corrections.
Consequently every formal series
\[
-\frac12\mu + \sum_{g\ge 1}\lambda_g\nu_g
\]
satisfies all homogeneous quadratic relations internal to this shell.
In particular, no quadratic identity confined to the first positive shell
can determine $D'_{1,1}$.
\end{proposition}

\begin{proof}
Each $\nu_g$ takes two inputs in $U_+$ and lands in $V_0$.
Therefore any insertion of one $\nu$-type map into another feeds a $V_0$
output into a slot that only accepts $U_+$, so the composition vanishes.
The map $\mu$ takes inputs in $U_+,U_+,\widetilde V_0$ and lands in $V_0$.
Inserting $\mu$ into itself again feeds a $V_0$ output into a slot that
requires either $U_+$ or $\widetilde V_0$, hence also vanishes.
The mixed compositions $\mu\circ \nu_g$ and $\nu_g\circ \mu$ vanish for the
same source-target reason.
Thus every bracket among these operations is zero.
\end{proof}

\begin{construction}[Partial low-mode quantum seed;
\ClaimStatusProvedHere]
\label{def:frontier-partial-low-mode-quantum-seed}
The previous theorem package canonically extracts from celestial BF a
minimal quantum seed
\[
\mathfrak S_{\mathrm{BF}}^{\mathrm{low}}(\Lambda)
=
\bigl(U_+\oplus V_0\oplus \widetilde V_0;
 m_2^{\mathrm{proj}},\ \mu,\ \Lambda(\hbar)\nu\bigr),
\]
where $m_2^{\mathrm{proj}}$ is the projected classical product,
\[
\mu(u,v,\widetilde\nu_0)=-\frac12\,\omega(u,v)\,\nu_0,
\qquad
\nu(u,v)=\omega(u,v)\,\nu_0,
\]
and
\[
\Lambda(\hbar)=\sum_{g\ge 1}\lambda_g\hbar^g.
\]
The actual celestial BF theory singles out
\[
\lambda_1 = D'_{1,1}h^\vee.
\]
This construction is the exact low-mode package of the first classical
correction and the first adjoint genus-$1$ correction.
\end{construction}

\begin{remark}[What the first positive shell does and does not know]
\label{rem:frontier-positive-shell-what-it-knows}
The first positive shell sees the determinant line
$\Lambda^2U_+\cong \mathbb C$ and nothing else.
This is already enough to prove formal rigidity of the shell and to reduce
all possible bilinear quantum corrections to a single scalar.
It is not enough to determine that scalar.
The coefficient $D'_{1,1}$ first appears as a parameter because every
quadratic relation internal to the determinant shell vanishes identically.
To constrain $D'_{1,1}$ one must enlarge the shell.
The smallest such enlargement is the axionic one.
\end{remark}

\subsection{The five-generator axionic shell and the first Jacobi equation}
\label{subsec:frontier-five-generator-shell}

The Kodaira--Spencer coupling supplies the unique scalar generator that can
interact nontrivially with the determinant shell at the lowest possible KK
weight:
\[
F:=\widetilde\beta_{0,0}.
\]
We now enlarge the first positive shell by adjoining $F$.
This is the first place where Jacobi can see the one-loop adjoint
coefficient $D'_{1,1}$.
The point is structural.
The determinant shell by itself is formally abelian, but the axionic sector
produces a linear $\widetilde\beta$-term in the $J\bar J$ OPE and a
second-order pole in the $jF$ OPE.
Together they force the first nontrivial Jacobi defect.

Let $\mathfrak g$ be a simple complex Lie algebra, and write
\[
J_a:=b^a_{1,0},
\qquad
\bar J_a:=b^a_{0,1},
\qquad
j_a:=b^a_{0,0},
\qquad
\widetilde j_a:=\widetilde b^a_{0,0},
\qquad
F:=\widetilde\beta_{0,0}.
\]
Let $\mathfrak S_5^{\partial}$ denote the free
$\mathbb C[\partial]$-module on these generators.
Define the linear projection
\[
\Pi_{\mathrm{lin}}\colon \text{full boundary algebra}
\longrightarrow \mathfrak S_5^{\partial}
\]
by retaining only linear terms in the displayed generators and discarding
all higher KK modes and all nonlinear composites.
For fields $a,b\in \mathfrak S_5^{\partial}$ define the
\emph{linearized $\lambda$-bracket}
\[
[a_{\lambda}b]_{\mathrm{lin}}:=\Pi_{\mathrm{lin}}([a_{\lambda}b]_{\mathrm{full}}).
\]

\begin{lemma}[Linearized five-generator brackets;
\ClaimStatusProvedHere]
\label{lem:frontier-five-generator-brackets}
The linearized $\lambda$-brackets on $\mathfrak S_5^{\partial}$ satisfy
\begin{align*}
[J_a{}_{\lambda}\bar J_b]_{\mathrm{lin}}
&=
D'_{1,1}\,h^\vee f_{ab}{}^c\,\lambda j_c
-\frac12 K_{ab}F,
\\
[j_a{}_{\lambda}F]_{\mathrm{lin}}
&=
2(\partial-\lambda)\widetilde j_a,
\\
[J_a{}_{\lambda}F]_{\mathrm{lin}}&=0,
\qquad
[\bar J_a{}_{\lambda}F]_{\mathrm{lin}}=0.
\end{align*}
Moreover the current algebra action survives linearly:
\begin{align*}
[j_a{}_{\lambda}j_b]_{\mathrm{lin}}&=f_{ab}{}^c j_c,
&[j_a{}_{\lambda}J_b]_{\mathrm{lin}}&=f_{ab}{}^c J_c,
\\
[j_a{}_{\lambda}\bar J_b]_{\mathrm{lin}}&=f_{ab}{}^c \bar J_c,
&[j_a{}_{\lambda}\widetilde j_b]_{\mathrm{lin}}&=f_{ab}{}^c\widetilde j_c.
\end{align*}
\end{lemma}

\begin{proof}
Zeng's tree-level $BB$ OPE contains the classical current term
\[
B^a(z)B^b(0)\sim \frac{f_{ab}{}^c}{z}B^c(0)
\]
and the higher $m_n$ corrections.
For $J_a=b^a_{1,0}$ and $\bar J_b=b^b_{0,1}$ the classical pole-$1$
$B$-term lands in mode $(1,1)$ and is therefore discarded by
$\Pi_{\mathrm{lin}}$.
The tree-level axionic correction from the Kodaira--Spencer coupling is
\[
J_a(z)\bar J_b(0)
\sim
-\frac12\frac{K_{ab}}{z}F
\]
because the $\widetilde\beta$ term in Zeng's formula has coefficient
$-(ps-qr)/(p+q+r+s)=-1/2$ on the mode pair $(1,0),(0,1)$.
The one-loop adjoint term contributes
\[
J_a(z)\bar J_b(0)
\sim
D'_{1,1}\frac{h^\vee f_{ab}{}^c}{z^2}j_c.
\]
Translating poles to $\lambda$-brackets gives the first formula.

For the second formula, set $(p,q,r,s)=(0,0,0,0)$ in the
$B\widetilde\beta$ coupling:
\[
j_a(z)F(0)
\sim
-\frac{2}{z^2}\widetilde j_a
+\frac{2}{z}\partial\widetilde j_a.
\]
Hence
\[
[j_a{}_{\lambda}F]_{\mathrm{lin}}
=
-2\lambda\widetilde j_a + 2\partial\widetilde j_a
=
2(\partial-\lambda)\widetilde j_a.
\]

For $J_a$ and $\bar J_a$ against $F$, Zeng's $B\widetilde\beta$ formula
lands in the modes $\widetilde B_{1,0}$ and $\widetilde B_{0,1}$,
respectively, together with their derivatives.
These are outside the five-generator linear shell, so their linearized
projections vanish.

The remaining current-action brackets are the classical current algebra
relations together with the classical $B\widetilde B$ OPE, all in mode
$(0,0)$, so they survive the linear projection exactly as stated.
\end{proof}

\begin{theorem}[First Jacobi equation on the five-generator shell;
\ClaimStatusProvedHere]
\label{thm:frontier-five-generator-jacobi}
Let
\[
\operatorname{Jac}_{\mathrm{lin}}(a,b,c)
:=[a_{\lambda}[b_{\mu}c]_{\mathrm{lin}}]_{\mathrm{lin}}
-[b_{\mu}[a_{\lambda}c]_{\mathrm{lin}}]_{\mathrm{lin}}
-[[a_{\lambda}b]_{\mathrm{lin},\lambda+\mu}c]_{\mathrm{lin}}
\]
denote the Jacobiator of the linearized $\lambda$-bracket.
Then for every pair of color labels $a,b$ one has
\[
\operatorname{Jac}_{\mathrm{lin}}(J_a,\bar J_b,F)
=
2D'_{1,1}\,h^\vee f_{ab}{}^c\,
\lambda(\lambda+\mu-\partial)\widetilde j_c.
\]
Equivalently,
\[
[[J_a{}_{\lambda}\bar J_b]_{\mathrm{lin},\lambda+\mu}F]_{\mathrm{lin}}
=
2D'_{1,1}\,h^\vee f_{ab}{}^c\,
\lambda(\partial-\lambda-\mu)\widetilde j_c.
\]
\end{theorem}

\begin{proof}
By
Lemma~\ref{lem:frontier-five-generator-brackets},
\([
\bar J_b{}_{\mu}F]_{\mathrm{lin}}=0=[J_a{}_{\lambda}F]_{\mathrm{lin}}
\)
so the first two terms in the Jacobiator vanish.
Thus
\[
\operatorname{Jac}_{\mathrm{lin}}(J_a,\bar J_b,F)
=
-[[J_a{}_{\lambda}\bar J_b]_{\mathrm{lin},\lambda+\mu}F]_{\mathrm{lin}}.
\]
Substitute the bracket from
Lemma~\ref{lem:frontier-five-generator-brackets}:
\[
[J_a{}_{\lambda}\bar J_b]_{\mathrm{lin}}
=
D'_{1,1}h^\vee f_{ab}{}^c\,\lambda j_c
-\frac12 K_{ab}F.
\]
The $F$-term contributes nothing inside the linearized shell.
For the $j_c$-term we obtain
\begin{align*}
[[J_a{}_{\lambda}\bar J_b]_{\mathrm{lin},\lambda+\mu}F]_{\mathrm{lin}}
&=
D'_{1,1}h^\vee f_{ab}{}^c\,\lambda
[j_c{}_{\lambda+\mu}F]_{\mathrm{lin}}
\\
&=
D'_{1,1}h^\vee f_{ab}{}^c\,\lambda\cdot
2(\partial-\lambda-\mu)\widetilde j_c.
\end{align*}
Negating this gives the displayed Jacobiator formula.
\end{proof}

\begin{corollary}[Strict five-generator closure forces vanishing of
$D'_{1,1}$; \ClaimStatusProvedHere]
\label{cor:frontier-strict-five-generator-closure-vanishing}
Assume $\mathfrak g$ is nonabelian.
If the linearized five-generator shell $\mathfrak S_5^{\partial}$ is
required to satisfy the Jacobi identity on its own, then necessarily
\[
D'_{1,1}=0.
\]
\end{corollary}

\begin{proof}
By
Theorem~\ref{thm:frontier-five-generator-jacobi},
Jacobi on the triple $(J_a,\bar J_b,F)$ forces
\[
2D'_{1,1}h^\vee f_{ab}{}^c\,
\lambda(\lambda+\mu-\partial)\widetilde j_c=0
\]
for all $a,b$.
For a nonabelian simple Lie algebra one has $h^\vee\neq 0$ and
$f_{ab}{}^c\not\equiv 0$, hence $D'_{1,1}=0$.
\end{proof}

\begin{definition}[Off-shell return term and axionic escape cocycle;
\ClaimStatusProvedHere]
\label{def:frontier-axionic-escape-cocycle}
Decompose the full bracket of $J_a$ and $\bar J_b$ as
\[
[J_a{}_{\lambda}\bar J_b]_{\mathrm{full}}
=
[J_a{}_{\lambda}\bar J_b]_{\mathrm{lin}}
+
[J_a{}_{\lambda}\bar J_b]_{\mathrm{off}},
\]
where the off-shell part contains every omitted contribution:
higher KK modes, nonlinear composites, and all fields outside the five
linear generators.
Define the projection
\[
\Pi_{\widetilde V_0}\colon
\text{full boundary algebra}
\longrightarrow
\mathbb C[\partial]\otimes (\mathfrak g\otimes \widetilde V_0)
\]
onto the $\widetilde j$-line, and set
\begin{align*}
\mathsf E_{ab}(\lambda,\mu)
:=
\Pi_{\widetilde V_0}\Bigl(&
[J_a{}_{\lambda}[\bar J_b{}_{\mu}F]_{\mathrm{full}}]_{\mathrm{full}}
-[\bar J_b{}_{\mu}[J_a{}_{\lambda}F]_{\mathrm{full}}]_{\mathrm{full}}
\\
&-[[J_a{}_{\lambda}\bar J_b]_{\mathrm{off},\lambda+\mu}F]_{\mathrm{full}}
\Bigr).
\end{align*}
We call $\mathsf E_{ab}(\lambda,\mu)$ the
\emph{axionic escape cocycle}.
It measures the precise amount by which omitted channels return to the
lowest $\widetilde j$-line in the Jacobi identity.
\end{definition}

\begin{corollary}[Exact first equation for $D'_{1,1}$;
\ClaimStatusProvedHere]
\label{cor:frontier-first-equation-for-D11}
In the full boundary algebra, the axionic escape cocycle satisfies
\[
\mathsf E_{ab}(\lambda,\mu)
=
2D'_{1,1}\,h^\vee f_{ab}{}^c\,
\lambda(\partial-\lambda-\mu)\widetilde j_c.
\]
Equivalently, $D'_{1,1}$ is exactly the coefficient of the universal tensor
\[
h^\vee f_{ab}{}^c\,\lambda(\partial-\lambda-\mu)\widetilde j_c
\]
in the $\widetilde j$-component of the omitted-sector Jacobi return.
\end{corollary}

\begin{proof}
Apply the full Jacobi identity to the triple $(J_a,\bar J_b,F)$ and split
$[J_a{}_{\lambda}\bar J_b]_{\mathrm{full}}$ into its linear and off-shell
parts.
Rearranging terms gives
\begin{align*}
\Pi_{\widetilde V_0}\Bigl(
&[J_a{}_{\lambda}[\bar J_b{}_{\mu}F]_{\mathrm{full}}]_{\mathrm{full}}
-[\bar J_b{}_{\mu}[J_a{}_{\lambda}F]_{\mathrm{full}}]_{\mathrm{full}}
\\
&-[[J_a{}_{\lambda}\bar J_b]_{\mathrm{off},\lambda+\mu}F]_{\mathrm{full}}
\Bigr)
=
\Pi_{\widetilde V_0}
[[J_a{}_{\lambda}\bar J_b]_{\mathrm{lin},\lambda+\mu}F]_{\mathrm{lin}}.
\end{align*}
The right-hand side is exactly the expression computed in
Theorem~\ref{thm:frontier-five-generator-jacobi}.
This is the displayed formula.
\end{proof}

\begin{remark}[The first genuine quantum equation]
\label{rem:frontier-first-genuine-quantum-equation}
Corollary~\ref{cor:frontier-first-equation-for-D11} is the first actual
equation for the one-loop coefficient $D'_{1,1}$.
It does not yet compute $D'_{1,1}$ numerically, but it does something more
structural: it identifies the exact tensor that must be produced by the
omitted channels.
Thus the problem of determining $D'_{1,1}$ has been reduced from a diffuse
regularization problem to a concrete return-map computation.
Any nonzero $D'_{1,1}$ must be carried by omitted sectors already visible in
Zeng's formulas, namely by some combination of:
\begin{itemize}
\item the off-shell $B[1,1]$ channel in the classical $BB$ OPE,
\item the axionic channels $\beta_{1,1}$ and $\partial\beta_{1,1}$,
\item the off-shell $\widetilde B_{1,0}$ and $\widetilde B_{0,1}$ channels in
  the $B\widetilde\beta$ OPE,
\item and the nonlinear composite channels generated by the higher
  $m_n$-terms.
\end{itemize}
In particular, the five-generator shell itself cannot support a nonzero
$D'_{1,1}$ internally; any such coefficient must be imported from omitted
channels and then returned to the $\widetilde j$-line by Jacobi.
\end{remark}

\begin{insight}[One mechanism, two frontiers]
\label{ins:frontier-one-mechanism-two-frontiers}
The higher-genus modular comparison problem and the five-generator celestial
BF problem are the same mathematical pattern at different resolutions.
In the global modular problem, once genus-$0$ normalization is fixed, the
only remaining ambiguity is a twisted cohomology class detected by a trace.
In the local celestial BF problem, once the determinant shell is fixed, the
only remaining ambiguity is the coefficient of a single universal tensor in
the omitted-sector Jacobi return.
In both cases the frontier is no longer ``mysterious higher structure.''
It is an explicit cohomology class.
\end{insight}

\begin{problem}[The next explicit computation]
\label{prob:frontier-compute-axionic-escape-cocycle}
Compute the axionic escape cocycle
$\mathsf E_{ab}(\lambda,\mu)$
from the first omitted packet visible in Zeng's formulas and determine the
resulting value of $D'_{1,1}$.
Equivalently, compute the $\widetilde j$-component of the omitted-sector
Jacobi return for the triple $(J_a,\bar J_b,F)$.
This is the smallest explicit renormalization problem whose solution would
produce the first fully determined genus-$1$ boundary correction in the
celestial BF example.
\end{problem}
